% manuscript_v5.tex
% Integrates: v3 (full manuscript), v4 (AUC gap decomposition),
%   step_16 (ML-BA predictor), step_17 (ML-BA external application),
%   step_18 (BA replacement experiment), single-feature BA analysis.
% Key narrative: cross-cultural gap dominates; BA method/availability is minor.

\documentclass[11pt,a4paper]{article}
\usepackage[margin=2.2cm]{geometry}
\usepackage[T1]{fontenc}
\usepackage[utf8]{inputenc}
\usepackage{booktabs}
\usepackage{longtable}
\usepackage{graphicx}
\usepackage{amsmath}
\usepackage{hyperref}
\usepackage{setspace}
\usepackage{microtype}
\usepackage{array}
\usepackage{xcolor}

\setstretch{1.08}
\setlength{\parskip}{4pt}
\setlength{\parindent}{0pt}

\hypersetup{
  colorlinks=true,
  linkcolor=blue,
  urlcolor=blue,
  pdftitle={KOA Risk Prediction: Internal Validation and Cross-Cultural Transportability},
  pdfauthor={BioAge-KOA-Prediction Team}
}

\title{Predicting Symptomatic Knee Osteoarthritis Risk from CHARLS Data:\\
\large Internal Validation, Biological Age Decomposition, and\\
Cross-Cultural Transportability Assessment on HRS and ELSA}
\author{BioAge-KOA-Prediction Team}
\date{July 2026}

\begin{document}
\maketitle

% ============================================================
\begin{abstract}
% ============================================================
\textbf{Objective.} To develop and internally validate a machine-learning model for
symptomatic knee osteoarthritis (KOA) risk prediction in a CHARLS-derived Chinese cohort,
to decompose the contribution of biological age (BA) to model performance, and to
characterise transportability to US (HRS) and UK (ELSA) aging cohorts.

\textbf{Methods.} A Random Forest model ($n=300$ trees) was trained on 17
sociodemographic and clinical variables (raw17) from CHARLS Dataset~A
($n=12{,}329$, KOA prevalence $13.3\%$). Internal validation used stratified 5-fold
out-of-fold estimation and an early 80/20 stratified unseen holdout. To isolate the
contribution of the pre-computed CHARLS ``Biological Age'' field --- whose derivation
method is undocumented --- we (i) computed transparent alternatives using the
Klemera--Doubal method (KDM-BA) and an adapted PhenoAge surrogate from 8 blood
biomarkers, (ii) replaced Sheet1 BA with each alternative inside raw17 while holding
feature count constant, and (iii) evaluated each BA variant as a single-feature KOA
predictor. To test whether a demographic-predicted BA could transport the aging signal
to cohorts lacking biomarkers, we trained ML regressors (Random Forest, XGBoost,
LightGBM ensemble) to predict both Sheet1 BA and KDM-BA from raw16 features alone
($R^2=0.64$ and $0.48$, respectively), then applied the predicted BA to HRS Wave~13
($n=20{,}605$) and ELSA Wave~8 ($n=8{,}416$) using harmonised raw16 feature sets.

\textbf{Results.} The primary model achieved ROC-AUC $0.847$ (internal OOF) / $0.907$
(unseen holdout). Replacing Sheet1 BA with KDM-BA degraded RF AUROC from $0.896$ to
$0.831$ ($\Delta=-0.065$); removing BA entirely produced $0.859$ ($\Delta=-0.037$).
As a single feature, all BA variants achieved only $0.50$--$0.54$ AUROC in CHARLS,
barely above chance. ML-predicted BA applied to external cohorts improved over raw16 by
only $+0.003$ to $+0.010$ AUROC. Restricting CHARLS to raw16 produced ROC-AUC $0.804$ /
$0.869$ --- a feature-subset penalty of $\Delta=-0.038$. Applying the same raw16 model to
ELSA and HRS yielded ROC-AUC $0.604$ for both cohorts --- a cross-cultural penalty of
$\Delta=-0.265$. The total AUC gap ($-0.303$) decomposes as $13\%$ feature-subset vs.\
$87\%$ cross-cultural. The RAS experiment recovered only $\Delta=+0.006$ AUC.

\textbf{Conclusions.} The biomarker--BA--KOA relationship is population-specific: BA is
a weak standalone KOA predictor ($r=0.02$--$0.04$) even in CHARLS, and demographic-predicted
BA adds negligible signal to external cohorts. The dominant driver of performance
degradation across populations is cross-cultural --- outcome definition heterogeneity,
diagnostic ascertainment, and population demographic differences --- not the absence of
Biological Age or the choice of BA computation method. These findings reframe the external
validation challenge: the primary obstacle is not a missing feature but a non-transferable
disease--aging relationship.

\textbf{Keywords:} symptomatic KOA, machine learning, biological age, KDM, cross-cultural
transportability, TRIPOD, calibration, DCA, HRS, ELSA
\end{abstract}

% ============================================================
\section{Introduction}
% ============================================================
Knee osteoarthritis is one of the leading causes of years lived with disability worldwide, with
age-standardised prevalence having risen steadily in recent decades\cite{Safiri2020}. In aging
populations, symptomatic KOA --- defined as clinically diagnosed knee osteoarthritis co-occurring with
knee pain --- imposes substantial mobility, quality-of-life, and health-care-utilization costs.
Risk-prediction models may support earlier stratification, but publication-quality evidence requires
more than high discrimination alone. Reporting should include calibration behaviour and threshold-level
decision consequences, not only ROC-AUC, and should follow standards such as TRIPOD\cite{Collins2015}
and PROBAST\cite{Wolff2019}.

A recent CHARLS-based report\cite{Fu2025} presented strong discrimination for symptomatic KOA and
highlighted biological age as a central predictor. However, two concerns arise when interpreting
this finding. First, the ``Biological Age'' field supplied in the CHARLS dataset is pre-computed
without documented methodology, raising questions about reproducibility and potential overfitting.
Second, translating a BA-centric prediction model to other populations assumes that the
biomarker--aging--disease relationship is biologically universal --- an assumption that has not been
tested across culturally distinct cohorts. This manuscript frames the project as an internal-validation
study with a structured decomposition of the BA contribution and a cross-cultural replication on the
RAND Health and Retirement Study (HRS) and the English Longitudinal Study of Ageing (ELSA) to
characterise the transportability of the aging--KOA signal across populations.

% ============================================================
\section{Methods}
% ============================================================

\subsection{Study Objective and Outcome}
The prediction target was symptomatic KOA in a CHARLS-derived cohort. The outcome was operationalised
from two CHARLS-derived fields as
\[
\text{KOA} = \mathbb{1}\bigl[\text{Arthritis} = \text{``yes''}\bigr]\ \wedge\
\mathbb{1}\bigl[\text{position\_knees} = \text{``yes''}\bigr],
\]
\noindent which matches the ``symptomatic KOA'' operationalisation used by Fu et al.\cite{Fu2025}
(physician-reported arthritis with knee-localised joint pain). Both source fields were removed from the
predictor set to prevent target leakage; rows with \texttt{position\_knees} missing were dropped because
the target itself could not be verified for those subjects.

\subsection{Study Cohort and Exclusion Cascade}
The analysis was built on Sheet1 of the CHARLS extract (15{,}545 baseline rows, 28 columns). Applying
the target-verification rule removed 3{,}216 rows with \texttt{position\_knees} missing, yielding the
analytical cohort ``Dataset A'' with 12{,}329 rows and 13.32\% symptomatic-KOA prevalence. Dataset A
was then partitioned once, before any model development, into an internal training partition
(80\%, $n=9{,}863$, prevalence 13.32\%) and an unseen holdout partition (20\%, $n=2{,}466$, prevalence
13.30\%) using a stratified split with \texttt{random\_state=42}. Table~\ref{tab:table1} summarises
the characteristics of Dataset A and its two partitions.

\begin{table}[h!]
\centering
\small
\caption{Cohort characteristics of the CHARLS-derived Dataset A and the early 80/20 partitions.
Values are counts (\%) or mean $\pm$ SD.}
\label{tab:table1}
\begin{tabular}{lccc}
\toprule
Characteristic & Full Dataset A & Internal train (80\%) & Unseen holdout (20\%) \\
\midrule
N & 12{,}329 & 9{,}863 & 2{,}466 \\
Symptomatic KOA, n (\%) & 1{,}642 (13.32) & 1{,}314 (13.32) & 328 (13.30) \\
Biological Age, mean $\pm$ SD & $58.70 \pm 9.79$ & $58.70 \pm 9.76$ & $58.70 \pm 9.90$ \\
Age bucket $\geq 60$, n (\%) & 3{,}357 (27.2) & 2{,}676 (27.1) & 681 (27.6) \\
Female, n (\%) & 6{,}443 (52.3) & 5{,}143 (52.1) & 1{,}300 (52.7) \\
BMI, mean $\pm$ SD & $23.65 \pm 3.84$ & $23.63 \pm 3.83$ & $23.73 \pm 3.88$ \\
BMI overweight, n (\%) & 10{,}122 (82.1) & 8{,}081 (81.9) & 2{,}041 (82.8) \\
BMI obese, n (\%) & 1{,}488 (12.1) & 1{,}192 (12.1) & 296 (12.0) \\
Married, n (\%) & 10{,}429 (84.6) & 8{,}343 (84.6) & 2{,}086 (84.6) \\
Education low/medium/high, \% & 28.1 / 61.9 / 10.0 & 27.9 / 62.0 / 10.1 & 28.8 / 61.4 / 9.7 \\
Residence urban, n (\%) & 7{,}954 (64.5) & 6{,}388 (64.8) & 1{,}566 (63.5) \\
Hypertension, n (\%) & 3{,}142 (25.5) & 2{,}503 (25.4) & 639 (25.9) \\
Dyslipidemia, n (\%) & 1{,}108 (9.0) & 882 (8.9) & 226 (9.2) \\
Diabetes, n (\%) & 726 (5.9) & 584 (5.9) & 142 (5.8) \\
Cancer, n (\%) & 93 (0.8) & 75 (0.8) & 18 (0.7) \\
CVD, n (\%) & 1{,}304 (10.6) & 1{,}016 (10.3) & 288 (11.7) \\
Smoking ever, n (\%) & 4{,}896 (39.7) & 3{,}897 (39.5) & 999 (40.5) \\
Drinking ever, n (\%) & 5{,}144 (41.7) & 4{,}156 (42.1) & 988 (40.1) \\
\bottomrule
\end{tabular}
\end{table}

\subsection{Feature Set (raw17)}
The pre-specified primary feature set, referred to throughout as \emph{raw17}, contains 17 variables
(Table~\ref{tab:raw17}). The subset excluding Biological Age, wave, and Time is labelled \emph{raw16}
and is used as the harmonised feature set for external validation.

\begin{table}[h!]
\centering
\small
\caption{The 17 variables in the raw17 feature set used by the primary model.}
\label{tab:raw17}
\begin{tabular}{llp{6.2cm}}
\toprule
Variable & Type & CHARLS-derived meaning \\
\midrule
wave & integer & Survey wave identifier \\
Time & integer & Survey year indicator \\
Gender & binary & Self-reported sex (1=male, 2=female) \\
Age\_New & ordinal (2) & Age bucket (45--59, $\geq 60$) \\
Marital & binary & Marital status (married vs other) \\
Education & ordinal (3) & Educational attainment (low/medium/high) \\
Residence & binary & Urban vs rural \\
Hypertension & binary & Self-reported physician-diagnosed \\
Dyslipidemia & binary & Self-reported physician-diagnosed \\
Diabetes & binary & Self-reported physician-diagnosed \\
Cancer & binary & Self-reported physician-diagnosed \\
CVD & binary & Self-reported cardiovascular disease \\
Smoke & binary & Ever vs never smoker \\
Drink & binary & Ever vs never drinker \\
BMI & continuous & Measured body-mass index (kg/m\textsuperscript{2}) \\
BMI\_New & ordinal (3) & BMI category (underweight/normal, overweight, obese) \\
Biological Age & continuous & CHARLS-supplied biological-age field \\
\bottomrule
\end{tabular}
\end{table}

\subsection{Missing Data Handling}
After the target-verification exclusion, the raw17 feature set and the outcome were complete for every
row in Dataset A (0\% missingness). As a pipeline safety net and to support sensitivity branches with
additional biomarker features, missing values were handled with train-fitted simple imputation (median
for numeric, most-frequent for categorical). Imputers and standardisation (\texttt{StandardScaler})
were fit on training data only and then applied to validation or unseen data.

\subsection{Early Unseen Holdout Strategy}
The early 80/20 stratified split was created before model development and persisted as CSV artefacts
(\texttt{dataset\_A\_internal\_train.csv}, \texttt{dataset\_A\_unseen\_test.csv}). Internal model
development and OOF estimation used only the internal partition; the unseen partition was reserved for
a single-shot external-style evaluation.

\subsection{Primary and Secondary Modeling Strategy}
The pre-specified primary model was a Random Forest (\texttt{RandomForestClassifier},
\texttt{n\_estimators=300}, \texttt{random\_state=42}, all other hyperparameters at scikit-learn
defaults). No grid search, Bayesian search, or Optuna-based tuning was performed on this configuration.
The BA replacement experiment (Section~\ref{sec:ba-replace}) additionally evaluated XGBoost and
LightGBM for robustness.

\subsection{Class-Imbalance Handling}
No resampling (SMOTE, random under-/over-sampling) and no \texttt{class\_weight} adjustment were
applied to the primary Random Forest. PR-AUC was pre-specified as the primary discrimination metric
alongside ROC-AUC, following standard recommendations for imbalanced binary prediction.

\subsection{Evaluation Framework}
Internal validation used stratified 5-fold out-of-fold (OOF) estimation on the internal training
partition, followed by re-evaluation on the unseen holdout. Reported metrics:
\begin{itemize}
\item discrimination: ROC-AUC, PR-AUC;
\item calibration: Brier score, ECE, MCE, Cox logistic calibration slope and intercept,
      decile reliability diagram;
\item decision relevance: Vickers net benefit at thresholds 0.05--0.60 and threshold-band
      summaries for 0.10--0.30.
\end{itemize}
Probability calibration was assessed for both uncalibrated (``raw'') and isotonic-calibrated
(\texttt{CalibratedClassifierCV}, \texttt{method=``isotonic''}, 3-fold inner CV) variants.

\subsection{Interpretability}
Model interpretation included SHAP global importance\cite{Lundberg2017}, SHAP interaction analysis,
and local waterfall plots for representative cases.

\subsection{Differences vs.\ Reference Study (Fu et al., 2025)}
Table~\ref{tab:paper-delta} summarises the key methodological differences between the reference
study and the present pipeline.

\begin{table}[h!]
\centering
\caption{Methodological delta versus reference study.}
\label{tab:paper-delta}
\begin{tabular}{p{3.4cm}p{5.1cm}p{5.1cm}}
\toprule
Domain & Reference paper & Current pipeline \\
\midrule
Cohort framing & Symptomatic KOA in CHARLS-derived sample &
  Symptomatic KOA in CHARLS-derived sample; updated internal protocol with early holdout \\
Missing data & Study-specific exclusions / imputation choices &
  Train-fitted \texttt{SimpleImputer} within each split \\
Split protocol & 70/30 emphasised & Early 80/20 unseen holdout + internal OOF validation \\
Class-balance handling & Model-specific handling &
  No resampling, no \texttt{class\_weight} adjustment \\
Primary metric & ROC-AUC-centred & PR-AUC reported alongside ROC-AUC \\
Primary model & XGBoost-centred & Random Forest + raw17 pre-specified \\
\bottomrule
\end{tabular}
\end{table}

% ============================================================
\subsection{Biological Age Computation: KDM-BA and PhenoAge}
% ============================================================
\label{sec:kdm-ba}

The CHARLS dataset contains a pre-computed ``Biological Age'' field whose derivation method is not
documented in the source publication\cite{Fu2025}. To test whether this feature is reproducible using
transparent methods, we computed two alternative biological age constructs from 8 blood biomarkers
available in CHARLS: platelet count, C-reactive protein (CRP), HbA1c, creatinine, blood urea nitrogen
(BUN), total cholesterol, triglycerides, and systolic blood pressure (SBP).

\textbf{KDM-BA.} The Klemera--Doubal Method\cite{Klemera2006} computes BA as a weighted combination
of biomarker regressions on chronological age. For each biomarker $x_i$, a linear regression
$x_i = a_i + b_i \cdot \text{CA} + \epsilon_i$ is fitted, yielding intercept $a_i$, slope $b_i$, and
residual standard error $\sigma_i$. The KDM biological age is:
\[
\text{BA}_{\text{KDM}} = \frac{\sum_i \frac{(x_i - a_i)\, b_i}{\sigma_i^2} + \frac{\text{CA}}{\sigma_{\text{CA}}^2}}
{\sum_i \frac{b_i^2}{\sigma_i^2} + \frac{1}{\sigma_{\text{CA}}^2}}
\]
We computed KDM-BA from both log-transformed (\texttt{BA\_KDM\_log}) and original-scale
(\texttt{BA\_KDM\_orig}) biomarker values. The ``Biological Age'' column in Sheet1 was used as the
chronological age proxy, as CHARLS does not provide a continuous age variable (Age\_New is a binary
$\geq 60$ indicator).

\textbf{PhenoAge\_Adapted.} Because canonical Levine PhenoAge\cite{Levine2018} requires biomarkers
not available in CHARLS (e.g., albumin, lymphocyte count), we constructed an adapted phenotypic-age
surrogate using winsorised z-scores of the same 8 biomarkers, weighted by clinical plausibility
(CRP: $+1.40$, HbA1c: $+1.20$, creatinine: $+1.00$, BUN: $+0.90$, SBP: $+1.00$,
total cholesterol: $+0.60$, triglycerides: $+0.60$, platelet: $-0.40$) and mapped to age units via
a scaling transformation. This score is explicitly non-canonical and is labelled \emph{PhenoAge\_Adapted}
throughout to avoid misinterpretation.

\textbf{Biological Illness Risk (BIR).} BIR was defined as $\text{BA}_{\text{KDM}} - \text{CA}$,
capturing the acceleration or deceleration of biological relative to chronological age.

% ============================================================
\subsection{Biological Age Replacement Experiment}
% ============================================================
\label{sec:ba-replace}

To determine whether the pre-computed Sheet1 BA is irreplaceable or whether any transparent
biomarker-derived BA achieves equivalent performance, we conducted a controlled replacement
experiment (Step 18). Six configurations were evaluated using the same cross-validation protocol
as the feature impact analysis (3 seeds $\times$ 5-fold StratifiedKFold, Random Forest with
StandardScaler):

\begin{enumerate}
  \item \textbf{raw17 (Sheet1 BA)} --- baseline, 17 features with the original BA field.
  \item \textbf{raw16\_only (no BA)} --- 16 features, BA removed entirely.
  \item \textbf{raw16 + KDM\_log} --- KDM-BA (log-transformed) replaces Sheet1 BA, 17 features.
  \item \textbf{raw16 + KDM\_orig} --- KDM-BA (original scale) replaces Sheet1 BA, 17 features.
  \item \textbf{raw16 + PhenoAge\_Adapted} --- adapted PhenoAge replaces Sheet1 BA, 17 features.
  \item \textbf{raw16 + PhenoAge\_Accel} --- PhenoAge acceleration replaces Sheet1 BA, 17 features.
\end{enumerate}

All configurations use the same $n=12{,}329$ CHARLS Dataset A and the same evaluation harness.
XGBoost (200 trees, max depth 6) and LightGBM (200 trees) were additionally evaluated for
robustness.

% ============================================================
\subsection{ML-Predicted Biological Age for External Transport}
% ============================================================
\label{sec:ml-ba}

Because ELSA and HRS lack the biomarkers needed to compute KDM-BA or PhenoAge directly, we tested
whether a demographic-predicted BA could transport the aging signal. Two ML regressor targets were
trained on CHARLS internal training data using only raw16 features (the sociodemographic variables
available in all three cohorts):

\begin{enumerate}
  \item \textbf{ML\_BA} --- predict the Sheet1 ``Biological Age'' column (CV $R^2=0.636$,
        holdout $R^2=0.633$, MAE $=4.7$ years).
  \item \textbf{ML\_KDM\_BA} --- predict the computed KDM-BA (CV $R^2=0.482$,
        holdout $R^2=0.501$, MAE $=6.0$ years).
\end{enumerate}

Each target was predicted by an ensemble of Random Forest, XGBoost, and LightGBM regressors
(300 trees each). The fitted models were then applied to the harmonised ELSA and HRS raw16
data to generate predicted BA columns, which were added to the raw16 feature set for KOA
prediction. Four scenarios per cohort were evaluated:

\begin{enumerate}
  \item raw16 only (no BA --- baseline)
  \item raw16 + ML\_BA
  \item raw16 + ML\_KDM\_BA
  \item raw16 + ML\_BA + ML\_KDM\_BA (ensemble BA --- addresses the ensemble prediction option)
\end{enumerate}

% ============================================================
\subsection{AUC Gap Decomposition Strategy}
% ============================================================
\label{sec:auc-decomp}

To disentangle the feature-subset penalty (attributable to BA absence) from the cross-cultural
penalty (attributable to population differences), we conducted the transportability analysis in
three stages:

\begin{enumerate}
  \item \textbf{CHARLS raw16 internal baseline.} The primary model was re-trained and
        evaluated on the CHARLS internal partition using only the features common to
        CHARLS, ELSA, and HRS (raw16). Because both training and evaluation occur within
        CHARLS, this isolates the feature-subset penalty with no cross-cultural confound.
  \item \textbf{External application.} The raw16 model trained on CHARLS was applied to
        ELSA Wave~8 and HRS Wave~13 without re-fitting. The gap between CHARLS raw16
        holdout performance and external cohort performance is the cross-cultural
        attenuation, net of the feature-subset penalty.
  \item \textbf{Reduced Aging Score (RAS) sensitivity.} To test whether a harmonised
        biological age proxy could recover the feature-subset gap, we constructed RAS
        from four biomarkers available in both CHARLS and ELSA (hs-CRP, HbA1c, systolic
        BP, total cholesterol). RAS was trained as a linear predictor of CHARLS KDM-BA
        ($R^2=0.417$) and added to the raw16 model for ELSA external validation.
\end{enumerate}

\subsection{Reduced Aging Score (RAS)}
\label{sec:ras}

RAS was constructed as a linear combination of four biomarkers:
\[
\text{RAS} = \beta_0 + \beta_1 \cdot \ln(\text{CRP}) + \beta_2 \cdot \ln(\text{HbA1c})
           + \beta_3 \cdot \ln(\text{TC}) + \beta_4 \cdot \text{SBP}
\]
The coefficients $\beta_{0..4}$ were estimated by ordinary least squares regression of
CHARLS KDM-BA on the four CHARLS biomarker equivalents
(\texttt{crp\_mg.L}, \texttt{Hb A1c}, \texttt{TC\_mg.d L}, \texttt{sbp.mean}).
$R^2$ on CHARLS data was $0.417$ ($n=12{,}329$), indicating that RAS captures
$41.7\%$ of KDM-BA's internal variance.

For ELSA, biomarkers were converted to CHARLS-equivalent units: hs-CRP (mg/L)
log-transformed directly; HbA1c converted from IFCC (mmol/mol) to DCCT (\%) via
$\text{HbA1c}\% = (\text{mmol/mol} / 10.929) + 2.15$, then log-transformed; total
cholesterol converted from mmol/L to mg/dL ($\times 38.67$), then log-transformed;
systolic BP used directly (mmHg). After excluding ELSA respondents with missing
biomarker codes ($-8$, $-2$, $-1$), $2{,}499$ of $8{,}416$ respondents had complete
biomarker data for RAS computation.

% ============================================================
\subsection{HRS and ELSA Cross-Cultural Replication}
% ============================================================

\subsubsection*{HRS Cohort and Data Source}
The Health and Retirement Study (HRS)\cite{HRS2023} is a US nationally representative longitudinal
panel of adults aged $\geq 50$. For this replication we used the RAND HRS Longitudinal File 2022
(Version 1)\cite{RAND2023}. Wave 13 (survey year 2016) was selected as the target wave.

\subsubsection*{HRS Eligibility and KOA Proxy Definition}
Inclusion criteria: respondents with Wave 13 interview data and age $\geq 45$ years
($n=20{,}640$ after age filter). A two-component KOA proxy was constructed:
(1) current physician-reported arthritis (\texttt{r13arthr} = 1) and
(2) knee-relevant functional impairment (\texttt{r13stoopa} = 1).
This proxy results in a prevalence of 33.9\% ($n_{\text{pos}}=6{,}976$). Rows with missing
arthritis response were excluded ($n=35$), yielding a final analytic sample of $n=20{,}605$.

\subsubsection*{ELSA Cohort and Data Source}
The English Longitudinal Study of Ageing (ELSA) is a UK nationally representative longitudinal
panel of adults aged $\geq 50$. We used the Gateway Harmonized ELSA file (Version H, 2002--2023),
distributed through the UK Data Service (SN 5050). Wave 8 (2016--2017) was the primary target wave;
Wave 9 (2018--2019) was retained as a temporal sensitivity analysis.

\subsubsection*{ELSA Eligibility and KOA Proxy Definition}
Inclusion criteria: Wave 8 interview data and age $\geq 45$ years ($n=8{,}416$ after age filter).
A two-component KOA proxy was constructed:
(1) ever-diagnosed arthritis (\texttt{r8arthre} = 1) and
(2) difficulty stooping/kneeling/crouching (\texttt{r8stoopa} = 1).
This proxy results in a prevalence of $24.9\%$ ($n_{\text{pos}}=2{,}097$).

\subsubsection*{External Feature Harmonisation (raw16)}
Both HRS and ELSA were harmonised to the raw16 feature set (raw17 minus Biological Age, wave, and
Time). HRS yielded 12 of 17 raw17 features (Biological Age, Residence, Dyslipidemia, and wave
excluded); ELSA yielded 14 of 17 (Biological Age, wave, and Time excluded; Residence and
Dyslipidemia available). The trained CHARLS model was re-fitted on the raw16 subset of the CHARLS
internal training partition and then applied to each external cohort.

% ============================================================
\section{Results}
% ============================================================

\subsection{Primary Model Performance (Random Forest + raw17)}
Table~\ref{tab:primary-metrics} summarises the matched internal OOF and unseen-holdout performance
for the pre-specified primary model.

\begin{table}[h!]
\centering
\small
\caption{Primary model (Random Forest + raw17) metrics.}
\label{tab:primary-metrics}
\begin{tabular}{llccccccc}
\toprule
Set & Variant & ROC-AUC & PR-AUC & Brier & ECE & Slope & Intercept & F1 \\
\midrule
Internal OOF & Raw & 0.847 & 0.688 & 0.070 & 0.045 & 0.989 & $-0.040$ & 0.649 \\
Internal OOF & Isotonic & 0.839 & 0.676 & 0.069 & 0.042 & 1.607 & 1.000 & 0.653 \\
Unseen holdout & Raw & 0.907 & 0.791 & 0.053 & 0.056 & 1.216 & 0.178 & 0.801 \\
Unseen holdout & Isotonic & 0.900 & 0.779 & 0.051 & 0.060 & 1.742 & 1.088 & 0.799 \\
\bottomrule
\end{tabular}
\end{table}

\subsection{Decision-Curve Analysis}
Across the 10\%--30\% threshold band, the primary model produced strictly positive mean net benefit
against both treat-all and treat-none. The derived net-reduction-per-100-patients
figure\cite{Vickers2016} was $56.0$ (internal OOF) and $63.8$ (unseen holdout) avoided unnecessary
interventions versus treat-all; for isotonic probabilities the corresponding values were $56.5$ and
$65.8$, with relative reductions of $65.2\%$ and $75.9\%$ against the treat-all baseline.

% ============================================================
\subsection{Biological Age Replacement Experiment}
% ============================================================
\label{sec:results-ba-replace}

Table~\ref{tab:ba-replace} presents the BA replacement experiment. Replacing Sheet1 BA with
any transparent alternative degraded performance; removing BA entirely performed better than
replacing it with KDM-BA.

\begin{table}[h!]
\centering
\small
\caption{Biological Age replacement experiment (Random Forest, 3 seeds $\times$ 5-fold CV,
$n=12{,}329$). All configurations preserve the same 17-feature count except raw16\_only.}
\label{tab:ba-replace}
\begin{tabular}{lccccc}
\toprule
Configuration & Features & ROC-AUC & PR-AUC & Brier & $\Delta$ from raw17 \\
\midrule
raw17 (Sheet1 BA)       & 17 & 0.896 & 0.776 & 0.057 & --- (baseline) \\
raw16\_only (no BA)     & 16 & 0.859 & 0.578 & 0.071 & $-0.037$ \\
raw16 + PhenoAge\_Adapted & 17 & 0.841 & 0.561 & 0.083 & $-0.055$ \\
raw16 + KDM\_orig       & 17 & 0.831 & 0.525 & 0.086 & $-0.065$ \\
raw16 + KDM\_log        & 17 & 0.829 & 0.513 & 0.087 & $-0.067$ \\
raw16 + PhenoAge\_Accel & 17 & 0.798 & 0.444 & 0.094 & $-0.098$ \\
\bottomrule
\end{tabular}
\end{table}

Three observations emerge. First, Sheet1 BA is the best BA variant in the full model
($0.896$), but its contribution is modest ($-0.037$ when removed). Second, replacing BA
with any transparent alternative \emph{actively hurts}: raw16 + KDM\_orig ($0.831$) is
worse than raw16\_no\_BA ($0.859$), indicating multicollinearity noise from inserting a
redundant biomarker-derived feature. Third, the maximum penalty from any BA manipulation
(method replacement or removal) is $0.067$ AUROC --- small relative to the $0.26$
cross-cultural gap reported below.

% ============================================================
\subsection{Single-Feature BA Predictive Value Across Cohorts}
% ============================================================
\label{sec:single-feature}

Table~\ref{tab:single-feature} presents each BA variant as a standalone KOA predictor
(univariate logistic regression). This directly addresses the question: does BA carry
KOA signal on its own, and does that signal transfer?

\begin{table}[h!]
\centering
\small
\caption{Single-feature BA $\rightarrow$ KOA AUROC across all cohorts. ``---'' indicates
the BA variant cannot be computed (missing biomarkers for KDM/PhenoAge; Sheet1 BA only
in CHARLS). ML-predicted BA is a demographic proxy (Step 16/17).}
\label{tab:single-feature}
\begin{tabular}{lccc}
\toprule
Feature & CHARLS ($n{=}12{,}329$) & ELSA ($n{=}8{,}416$) & HRS ($n{=}20{,}605$) \\
\midrule
Biological Age (Sheet1)        & 0.536 & --- & --- \\
KDM-BA (log)                   & 0.516 & --- & --- \\
KDM-BA (orig)                  & 0.519 & --- & --- \\
PhenoAge\_Adapted              & 0.530 & --- & --- \\
PhenoAge\_Accel                & 0.504 & --- & --- \\
BIR\_log                       & 0.535 & --- & --- \\
ML\_BA (predicted)             & ---   & 0.556 & 0.629 \\
ML\_KDM\_BA (predicted)        & ---   & 0.589 & 0.629 \\
Age\_New (no BA, baseline)     & ---   & 0.553 & 0.599 \\
\bottomrule
\end{tabular}
\end{table}

All BA variants in CHARLS achieve only $0.50$--$0.54$ AUROC --- barely above chance.
The best (Sheet1 BA at $0.536$) has a KOA correlation of $r=0.038$.
This refutes the interpretation that BA is a strong standalone KOA predictor: its apparent importance in the full raw17 model (SHAP $=0.439$)
 reflects conditional interaction effects, not marginal predictive power.

In HRS, ML-predicted BA achieves $0.629$ --- but Age\_New alone achieves $0.599$ (HRS has
$33.9\%$ KOA prevalence vs.\ CHARLS's $13.3\%$). The BA increment over age is only
$+0.030$. In ELSA, ML\_BA ($0.556$) is nearly identical to Age\_New ($0.553$).
ML\_KDM\_BA consistently adds more over age than ML\_BA ($+0.036$ vs.\ $+0.003$ in ELSA;
$+0.030$ vs.\ $+0.030$ in HRS), suggesting that the biomarker-derived KDM target encodes
more KOA-relevant signal than the undocumented Sheet1 BA formula.

% ============================================================
\subsection{ML-Predicted BA Applied to External Cohorts}
% ============================================================
\label{sec:results-ml-ba-external}

Table~\ref{tab:ml-ba-external} presents the effect of adding ML-predicted BA to the
raw16 model in ELSA and HRS.

\begin{table}[h!]
\centering
\small
\caption{ML-predicted BA applied to external cohorts. All scenarios use the CHARLS-trained
Random Forest re-fitted on raw16. ML\_BA and ML\_KDM\_BA are demographic-predicted BA from
Step 16. The ensemble scenario (raw16 + ML\_BA + ML\_KDM\_BA) addresses the ensemble BA
prediction option.}
\label{tab:ml-ba-external}
\begin{tabular}{lcccccc}
\toprule
Cohort & Scenario & $n$ & ROC-AUC & PR-AUC & Brier & F1 \\
\midrule
ELSA & raw16               & 8{,}416  & 0.604 & 0.297 & 0.201 & 0.427 \\
ELSA & raw16 + ML\_BA       & 8{,}416  & 0.602 & 0.307 & 0.191 & 0.422 \\
ELSA & raw16 + ML\_KDM\_BA  & 8{,}416  & 0.597 & 0.300 & 0.192 & 0.421 \\
ELSA & raw16 + ML\_BA + ML\_KDM\_BA & 8{,}416 & 0.613 & 0.319 & 0.187 & 0.425 \\
\midrule
HRS  & raw16               & 20{,}605 & 0.604 & 0.399 & 0.258 & 0.525 \\
HRS  & raw16 + ML\_BA       & 20{,}605 & 0.607 & 0.410 & 0.250 & 0.526 \\
HRS  & raw16 + ML\_KDM\_BA  & 20{,}605 & 0.610 & 0.412 & 0.248 & 0.525 \\
HRS  & raw16 + ML\_BA + ML\_KDM\_BA & 20{,}605 & 0.607 & 0.412 & 0.247 & 0.523 \\
\bottomrule
\end{tabular}
\end{table}

The ensemble BA scenario (raw16 + ML\_BA + ML\_KDM\_BA) achieved the highest AUROC in
ELSA ($0.613$, $+0.009$ over raw16) but not in HRS ($0.607$). ML\_KDM\_BA alone was
best in HRS ($0.610$, $+0.006$). All gains are within $\pm 0.01$ of raw16 ---
negligible relative to the $0.26$ cross-cultural gap.

% ============================================================
\subsection{AUC Gap Decomposition: Feature-Subset vs.\ Cross-Cultural}
% ============================================================

Table~\ref{tab:auc-decomp} presents the structured decomposition of the
$0.303$ ROC-AUC gap between the CHARLS raw17 holdout and the external cohorts.

\begin{table}[h!]
\centering
\small
\caption{AUC gap decomposition across CHARLS raw17, CHARLS raw16, ELSA, and HRS.
All values are raw (uncalibrated) ROC-AUC except where noted.}
\label{tab:auc-decomp}
\begin{tabular}{lcccccc}
\toprule
Stage & Cohort & Features & $n$ & KOA\% & ROC-AUC & $\Delta$ from baseline \\
\midrule
A  & CHARLS raw17 Holdout  & 17 & 2{,}466  & 13.3 & 0.907 & --- (baseline) \\
B  & CHARLS raw16 Holdout  & 15 & 2{,}466  & 13.3 & 0.869 & $-0.038$ (feature-subset) \\
C  & ELSA W8 raw16         & 14 & 8{,}416  & 24.9 & 0.604 & $-0.303$ (total) \\
C' & HRS W13 raw16         & 12 & 20{,}605 & 33.9 & 0.604 & $-0.303$ (total) \\
\cmidrule{6-7}
   & \textbf{Feature-subset component}  & & & & & $\mathbf{-0.038}$ (13\%) \\
   & \textbf{Cross-cultural component}  & & & & & $\mathbf{-0.265}$ (87\%) \\
\bottomrule
\end{tabular}

\vspace{4pt}
\footnotesize{Feature-subset penalty = B $-$ A. Cross-cultural penalty = C $-$ B.
Percentages are of the total A$\rightarrow$C gap ($0.303$).}
\end{table}

The feature-subset penalty (dropping Biological Age, wave, and Time from raw17)
accounts for $13\%$ of the total AUC loss. The remaining $87\%$ is cross-cultural
attenuation occurring between the CHARLS and external cohorts.

Notably, the ELSA and HRS raw16 ROC-AUC values are nearly identical ($0.604$)
despite ELSA having two more harmonised features (Residence and Dyslipidemia), using
a different KOA proxy (ever-diagnosed vs.\ current-wave arthritis), and drawing from
a different population (UK vs.\ US). This convergence supports a stable transportability
ceiling around $\text{AUROC} \approx 0.60$ for raw16 without Biological Age across
culturally distinct aging cohorts.

\subsection{Reduced Aging Score (RAS) Experiment}

Table~\ref{tab:ras} presents the effect of adding RAS to the ELSA raw16 model.

\begin{table}[h!]
\centering
\small
\caption{Effect of the Reduced Aging Score on ELSA external validation.
RAS captures 41.7\% of KDM-BA variance in CHARLS.}
\label{tab:ras}
\begin{tabular}{lccc}
\toprule
Model & ROC-AUC & PR-AUC & Brier \\
\midrule
ELSA raw16              & 0.6050 & 0.297 & 0.200 \\
ELSA raw16 + RAS        & 0.6111 & 0.316 & 0.188 \\
\cmidrule{2-4}
$\Delta$ (RAS gain)     & \textbf{+0.0061} & +0.019 & $-0.013$ \\
\bottomrule
\end{tabular}
\end{table}

Adding RAS --- the best biological age proxy constructible from publicly available
ELSA biomarkers --- recovers only $0.006$ ROC-AUC, or approximately $2\%$ of the
total cross-cultural gap.

\subsection{ELSA Wave 9 Sensitivity Analysis}
A temporal sensitivity analysis using ELSA Wave 9 (2018--2019, $n=8{,}683$, KOA proxy
prevalence $23.3\%$) produced consistent results: ROC-AUC $0.596$ (raw), calibration
slope $0.187$. After split-calibrate-evaluate isotonic recalibration: ROC-AUC $0.615$,
slope $1.225$, ECE $0.009$. Wave 9 lacks a harmonized BMI variable, so BMI was fully
imputed from CHARLS training data. The results confirm temporal robustness.

% ============================================================
\section{Discussion}
% ============================================================

\paragraph{Internal validation.}
The primary raw17 model achieved ROC-AUC 0.847 (OOF) / 0.907 (holdout) and PR-AUC 0.688 / 0.791,
substantially exceeding the prevalence baseline. Strictly positive DCA net benefit in the 10\%--30\%
threshold band provides a grounded internal-evidence basis for risk-stratification decision support.

\paragraph{The BA contribution is modest and method-dependent.}
The BA replacement experiment (Table~\ref{tab:ba-replace}) reveals that Sheet1 BA contributes
$0.037$ AUROC to the raw17 model. Replacing it with transparent KDM-BA or PhenoAge degrades
performance further ($-0.055$ to $-0.098$), not because these methods are inferior BA estimators,
but because inserting a redundant biomarker-derived feature into raw16 introduces
multicollinearity noise. The single-feature analysis (Table~\ref{tab:single-feature})
confirms that all BA variants --- including Sheet1's --- are weak standalone KOA predictors
($0.50$--$0.54$ AUROC, $r=0.02$--$0.04$). BA's apparent importance in the full model
(SHAP $= 0.439$) reflects conditional interaction effects, not marginal predictive value.

\paragraph{The biomarker--BA--KOA relationship is population-specific.}
The central hypothesis tested in this study is that the aging--disease relationship
established in CHARLS does not transfer to other populations. Five lines of evidence
converge:

\begin{enumerate}
  \item \textbf{BA is weak even in CHARLS.} All BA variants achieve only $0.50$--$0.54$
        AUROC as single-feature KOA predictors ($r=0.02$--$0.04$). If BA were a strong
        biological signal for KOA, one would expect higher standalone discrimination.
  \item \textbf{Removing BA costs only $0.037$ AUROC.} The BA removal penalty is a small
        fraction of the total $0.303$ cross-cohort gap.
  \item \textbf{ML-predicted BA adds negligible signal externally.} Demographic-predicted
        BA (ML\_BA, ML\_KDM\_BA) improves raw16 by $+0.003$ to $+0.010$ AUROC in ELSA/HRS.
        The ensemble BA scenario adds $+0.009$. These gains are at least $25\times$ smaller
        than the cross-cultural gap.
  \item \textbf{ELSA--HRS convergence despite differences.} ELSA and HRS produce
        near-identical ROC-AUC ($0.604$) despite different feature availability (14 vs.\
        12 variables), KOA proxy definitions (ever-diagnosed vs.\ current-wave arthritis),
        and populations (UK vs.\ US). This convergence is not compatible with a BA-centric
        narrative: if missing BA were the primary driver, cohorts with different feature
        completeness should perform differently.
  \item \textbf{RAS recovers negligible signal.} The best biological age proxy
        constructible from ELSA biomarkers recovers $+0.006$ ROC-AUC --- less than $2\%$
        of the cross-cultural gap.
\end{enumerate}

\paragraph{AUC gap decomposition: cross-cultural dominates.}
The structured decomposition (Table~\ref{tab:auc-decomp}) shows that $87\%$ of the
total AUC gap is cross-cultural, not feature-subset. The complete evidence chain is:

\begin{table}[h!]
\centering
\small
\caption{Evidence chain: BA method choice and availability produce small penalties;
the cross-cultural cliff dominates.}
\label{tab:chain}
\begin{tabular}{llcc}
\toprule
Stage & Description & ROC-AUC & $\Delta$ \\
\midrule
1 & raw17 + Sheet1 BA (CHARLS)      & 0.896 & --- \\
2 & raw16 + KDM-BA (CHARLS)         & 0.831 & $-0.065$ (BA method) \\
3 & raw16 no BA (CHARLS)           & 0.859 & $-0.037$ (BA removal) \\
4 & raw16 no BA (ELSA)             & 0.604 & $-0.292$ (cross-cultural) \\
4 & raw16 no BA (HRS)              & 0.604 & $-0.292$ (cross-cultural) \\
5 & raw16 + ML-BA ensemble (ELSA)  & 0.613 & $+0.009$ (ML-BA recovers) \\
5 & raw16 + ML-KDM-BA (HRS)        & 0.610 & $+0.006$ (ML-BA recovers) \\
\bottomrule
\end{tabular}
\end{table}

No manipulation of BA --- using the original formula, replacing it with KDM-BA,
removing it entirely, or predicting it from demographics --- can account for the
$0.26$ cross-cultural gap. The BA-related penalties ($0.037$--$0.067$) are
$4$--$7\times$ smaller than the cross-cultural penalty ($0.265$). Even if full
biomarker-derived BA were available in ELSA/HRS, the expected recovery is bounded
by what BA contributes in CHARLS ($\sim 0.04$ AUROC), which is $<15\%$ of the gap.

\paragraph{Comparison with the reference study.}
The reference study by Fu et al.\cite{Fu2025} reported XGBoost-centred discrimination with
biological age as the dominant predictor (SHAP $> 0.6$), using a 70/30 split. The present
pipeline replicates the same outcome operationalisation and confirms that the raw17 feature
set supports strong discrimination. However, our analysis adds three findings that qualify
the original interpretation: (1) the pre-computed BA field has undocumented methodology
and is only moderately better than transparent KDM-BA ($0.896$ vs.\ $0.831$);
(2) BA is a weak standalone predictor ($0.536$ AUROC, $r=0.038$) whose apparent importance
reflects interaction effects; and (3) the BA--KOA relationship does not transport across
populations.

\paragraph{Strengths.}
Strengths include the large population-based cohort, transparent leakage controls,
pre-specification of the primary model, multi-dimensional evaluation (discrimination,
calibration, clinical utility, interpretability), adherence to TRIPOD and PROBAST,
the BA replacement experiment directly testing whether Sheet1 BA is reproducible,
the ML-BA transport experiment, and the structured AUC gap decomposition.

% ============================================================
\section{Limitations}
% ============================================================

\begin{itemize}

\item \textbf{Outcome proxy heterogeneity is the primary limitation.} The AUC gap
      decomposition shows that $87\%$ of performance loss is cross-cultural, with
      outcome definition mismatch (CHARLS knee-localised pain vs.\ HRS/ELSA stooping
      difficulty) likely the largest single contributor. Until a harmonised KOA
      definition is applied across all three cohorts, the true transportability of
      the raw17 model cannot be estimated.

\item \textbf{Sheet1 Biological Age is undocumented.} The pre-computed BA field in
      the CHARLS dataset lacks published methodology. Our analysis shows its residual
      (after removing age-group variance) correlates strongly with biomarkers
      (SBP $r=0.41$, creatinine $r=0.32$, BUN $r=0.32$), suggesting it was derived
      from the same biomarkers used in our KDM-BA computation. However, the exact
      formula --- and therefore its susceptibility to overfitting or information
      leakage --- cannot be verified.

\item \textbf{ML-predicted BA is a demographic proxy, not real BA.} The ML-BA models
      predict BA from sociodemographic features ($R^2=0.64$ for Sheet1 BA, $0.48$
      for KDM-BA). They capture only the demographic-BA relationship, not the
      biomarker-BA relationship. A reviewer could argue that real biomarker-derived
      BA in ELSA/HRS would perform better. However, given that BA contributes only
      $0.037$ AUROC in CHARLS and is a weak standalone predictor ($0.536$), the
      expected gain is bounded and cannot account for the $0.26$ gap.

\item \textbf{ELSA and HRS lack full biomarker panels.} KDM-BA cannot be computed in
      ELSA (missing creatinine, BUN, platelets) or HRS (requires VBS linkage). RAS,
      based on 4 of 8 KDM biomarkers, captures only $42\%$ of KDM-BA variance. We
      acknowledge that a cohort with complete biomarkers and the same KOA definition
      as CHARLS would provide a fairer test. However, the near-identical ELSA/HRS
      performance ($0.604$ each) despite different feature availability suggests
      feature completeness is not the binding constraint.

\item \textbf{Cross-sectional outcome.} The KOA label is cross-sectionally derived
      and cannot distinguish incident from prevalent cases.

\item \textbf{Self-reported comorbidities.} Hypertension, diabetes, cancer, CVD,
      dyslipidemia, and lifestyle variables are self-reported in CHARLS, inheriting
      recall bias and diagnostic ascertainment heterogeneity.

\item \textbf{Bootstrap confidence intervals not yet tabulated.} Point estimates and
      reliability diagrams are the current evidence basis for calibration; bootstrap
      CIs for discrimination and calibration metrics will be added in the next
      revision (TRIPOD item 16).

\end{itemize}

% ============================================================
\section{Conclusion}
% ============================================================

This study demonstrates a methodologically rigorous internal-validation pipeline for
symptomatic KOA prediction in a large Chinese nationally-representative cohort. The
pre-specified primary model (Random Forest + raw17) achieves strong discrimination
(ROC-AUC $0.907$ on the unseen holdout), acceptable calibration, and positive DCA net
benefit in the clinically relevant 10\%--30\% risk threshold band.

A structured decomposition of the biological age contribution reveals that BA ---
whether pre-computed by an undocumented formula, transparently computed via KDM, or
demographic-predicted via ML --- is a weak standalone KOA predictor ($0.50$--$0.54$
AUROC) whose contribution to the full model is modest ($0.037$ AUROC) and whose
relationship with KOA does not transfer across populations. ML-predicted BA applied
to external cohorts adds $+0.003$ to $+0.010$ AUROC; an ensemble of two ML-BA
variants adds $+0.009$. All BA-related effects combined are bounded at $\sim 0.067$
AUROC.

The cross-cultural gap ($0.265$ AUROC, $87\%$ of the total) is the dominant barrier
to transportability. The near-identical performance of the raw16 model on ELSA and
HRS ($0.604$) despite different feature availability, KOA proxy definitions, and
populations establishes a conservative transportability ceiling that is not explained
by BA absence. These findings reframe the external validation challenge: the primary
obstacle is not a missing feature but a non-transferable disease--aging relationship.
Future work should prioritise (i) harmonised KOA outcome definitions across cohorts,
(ii) prospective validation in cohorts where the same KOA definition and full
biomarker panels are available, and (iii) investigation of population-specific
effect modifiers that drive the cross-cultural gap.

% ============================================================
\section*{Data Availability}
% ============================================================
The CHARLS development dataset (``Raw Data of Biological Age'' derived from CHARLS) is publicly
available under a CC BY 4.0 licence at Mendeley Data:
\url{https://data.mendeley.com/datasets/3rv7mf5pv9/1} (Fu F, Mendeley Data V1, 2025;
doi:10.17632/3rv7mf5pv9.1). The RAND HRS Longitudinal File is publicly available at
\url{https://hrsdata.isr.umich.edu/data-products/rand} upon free registration. The English
Longitudinal Study of Ageing (ELSA) is available from the UK Data Service (SN 5050) at
\url{https://beta.ukdataservice.ac.uk/datacatalogue/studies/study?id=5050} under an End User
Licence. The complete analysis pipeline code is available at the project repository.

% ============================================================
\begin{thebibliography}{99}
% ============================================================

\bibitem{Fu2025} Fu F, Dong L, Lian J, et al. Biological age threshold is associated with
symptomatic knee osteoarthritis risk in Chinese adults: insights from machine learning analysis
of a national cohort. \textit{PLOS One}. 2025;20(12):e0335250.

\bibitem{Collins2015} Collins GS, Reitsma JB, Altman DG, Moons KGM. Transparent reporting of a
multivariable prediction model for individual prognosis or diagnosis (TRIPOD): the TRIPOD statement.
\textit{BMJ}. 2015;350:g7594.

\bibitem{Wolff2019} Wolff RF, Moons KGM, Riley RD, et al. PROBAST: a tool to assess the risk of
bias and applicability of prediction model studies. \textit{Ann Intern Med}. 2019;170(1):51--58.

\bibitem{Safiri2020} Safiri S, Kolahi AA, Smith E, et al. Global, regional and national burden of
osteoarthritis 1990--2017: a systematic analysis of the Global Burden of Disease Study 2017.
\textit{Ann Rheum Dis}. 2020;79(6):819--828.

\bibitem{Vickers2006} Vickers AJ, Elkin EB. Decision curve analysis: a novel method for evaluating
prediction models. \textit{Med Decis Making}. 2006;26(6):565--574.

\bibitem{Vickers2016} Vickers AJ, van Calster B, Steyerberg EW. Net benefit approaches to the
evaluation of prediction models, molecular markers, and diagnostic tests. \textit{BMJ}. 2016;352:i6.

\bibitem{VanCalster2019} Van Calster B, McLernon DJ, van Smeden M, Wynants L, Steyerberg EW.
Calibration: the Achilles heel of predictive analytics. \textit{BMC Med}. 2019;17(1):230.

\bibitem{NiculescuMizil2005} Niculescu-Mizil A, Caruana R. Predicting good probabilities with
supervised learning. In: \textit{Proceedings of the 22nd International Conference on Machine
Learning (ICML)}. 2005:625--632.

\bibitem{Steyerberg2019} Steyerberg EW. \textit{Clinical Prediction Models: A Practical Approach
to Development, Validation, and Updating}. 2nd ed. Springer; 2019.

\bibitem{Lundberg2017} Lundberg SM, Lee SI. A unified approach to interpreting model predictions.
In: \textit{Advances in Neural Information Processing Systems}. 2017.

\bibitem{Klemera2006} Klemera P, Doubal S. A new approach to the concept and computation of
biological age. \textit{Mechanisms of Ageing and Development}. 2006;127(3):240--248.

\bibitem{Levine2018} Levine ME, Lu AT, Quach A, et al. An epigenetic biomarker of aging for
lifespan and healthspan. \textit{Aging}. 2018;10(4):573--591.

\bibitem{HRS2023} Health and Retirement Study. \textit{Public Use Dataset}. Produced and distributed
by the University of Michigan with funding from the National Institute on Aging (grant number
NIA U01AG009740). Ann Arbor, MI; 2023.

\bibitem{RAND2023} RAND HRS Longitudinal File 2022 (V1). Santa Monica, CA: RAND Corporation; 2023.
\url{https://hrsdata.isr.umich.edu/data-products/rand}.

\bibitem{ELSA2025} Banks J, Batty GD, Breedvelt J, et al. English Longitudinal Study of Ageing:
Waves 0--11, 1998--2024. [data collection]. 50th Edition. UK Data Service. SN 5050, 2025.
doi:10.5255/UKDA-SN-5050-37.

\end{thebibliography}

\end{document}